\documentclass[10pt,a4paper]{amsart}
\usepackage[margin=19mm]{geometry}
\usepackage{amsmath,amssymb,mathtools}
\usepackage[scr=rsfs,cal=euler]{mathalfa}
\usepackage{xfrac}
\usepackage[unicode,hidelinks,bookmarks=false]{hyperref}
\newcommand{\ie}{i.e.\ }
\newcommand{\cf}{cf.\ }
\newcommand{\resp}{respectively\ }
\newcommand{\Subsection}{\subsection}
\providecommand{\nobreakdash}{\nobreak\hbox{-}}
\setlength{\emergencystretch}{2em}
\multlinegap0pt
\usepackage{bm}
\usepackage{bbm}
\usepackage{dsfont}
\usepackage{xspace}
\usepackage{cancel}
\usepackage{mathtools}
\usepackage[shortlabels]{enumitem}
\usepackage{amsthm}
\makeatletter
\@ifundefined{definition}{\newtheorem{definition}{Definition}}{}
\@ifundefined{claim}{\newtheorem{claim}[definition]{Claim}}{}
\@ifundefined{fact}{\newtheorem{fact}[definition]{Fact}}{}
\@ifundefined{proposition}{\newtheorem{proposition}[definition]{Proposition}}{}
\makeatother
\title[Spectral generalized Tur\'{a}n problems]{Spectral generalized Tur\'{a}n problems}
\author{Xizhi Liu}
\date{Source: arXiv:2507.21689v1 (CC BY 4.0)}
\begin{document}
\maketitle

\noindent\emph{Source and licence.} Spectral generalized Tur\'{a}n problems. Version arXiv:2507.21689v1 (CC BY 4.0), distributed under the Creative Commons Attribution 4.0 International licence, as recorded in the arXiv licence metadata for this exact version and shown on the abstract page https://arxiv.org/abs/2507.21689v1. Full rights notice: licenses/arxiv-2507.21689-CC-BY-4.0.txt.\par\smallskip

\noindent\textit{Editorial scope.} Counted unit: the Proposition whose source label is \texttt{PROP:Lagrangian-Entropy} in the article (source0.tex lines 1220--1236, source label \texttt{PROP:Lagrangian-Entropy}), delivered together with the article's own complete original proof of it (source0.tex lines 1237--1473, one \texttt{proof} environment of 237 lines). No mathematics is written, repaired or re-derived: statement and proof are the article's own text line for line. Required context retained uncounted: FACT:nonnegative-optimal-vector (lines 419--421); prop:random-embedding-b (lines 1200--1203). Every definition, display and label of those lines is retained unchanged. Omitted from this package and not redistributed: 1--418, 422--1199, 1204--1219, 1474--1650. Nothing inside a retained excerpt is omitted or altered: every retained body is delivered exactly as the article's lines are written, so no omission receipt of any kind accompanies this package. Citations: the retained excerpts carry no citation command at all, so no bibliography entry of the article is transcribed and no reference is redistributed. Reference adaptation: none was needed and none was made; every reference inside the delivered unit resolves inside it, so no unresolved reference or citation is produced. Printed slips kept literal: no printed word, symbol, hypothesis or number of the counted statement or of its proof was altered. Layout and class: the article's own class is \texttt{article} with packages including epsf, epsfig, amsfonts, amsgen, amsmath, amstext, amsbsy, amsopn, amsthm, cases, listings, color, ebezier, eepic; the wrapper is the lane's pinned \texttt{amsart} editorial wrapper at 10pt on A4, which restates the theorem-like environments the retained text needs and adds the article's own macro definitions that the retained text needs (none), with extra wrapper packages (bm, bbm, dsfont, xspace, cancel, mathtools, enumitem). The delivered numbering is the unit's own. Assets: no retained span contains a figure environment, a graphics-inclusion command, a PDF-page inclusion, a table or a TikZ picture; no image file is copied into this package, the package declares no asset dependency and no image file is redistributed with it. Licence: version arXiv:2507.21689v1 of this article is distributed under the Creative Commons Attribution 4.0 International licence, as recorded in the arXiv licence metadata for this exact version and shown on the abstract page https://arxiv.org/abs/2507.21689v1. A full rights notice is delivered at licenses/arxiv-2507.21689-CC-BY-4.0.txt. Characters: every retained excerpt is pure ASCII, so the delivered engine, XeLaTeX with its default Latin Modern text font, reports no missing character.
\begin{fact}\label{FACT:nonnegative-optimal-vector}
    The set $\mathrm{OPT}_{\alpha, Q}(\mathcal{H})$ contains a nonnegative vector. 
\end{fact}

\begin{enumerate}[label=(\roman*)]
    \item\label{prop:random-embedding-a} $(X_1, \ldots, X_{r})$ is always an element in $\mathrm{Inj}(Q,\mathcal{H})$, and 
    \item\label{prop:random-embedding-b} $(X_1, \ldots, X_{r})$ is symmetric with respect to $\mathrm{Aut}(Q)$, that is, $(X_1, \ldots, X_{q})$ is the same as $(X_{\sigma(1)}, \ldots, X_{\sigma(q)})$ for every $\sigma \in \mathrm{Aut}(Q)$. 
\end{enumerate}

\begin{proposition}\label{PROP:Lagrangian-Entropy}
    % We have 
    % \begin{align}\label{equ:entropy-vs-Lagrangian-a}
    %     \eta_{\alpha,Q}(\mathcal{H})
    %     \ge \lambda_{\alpha,Q}(\mathcal{H}).
    % \end{align}
    % Moreover, if $Q$ is vertex-transitive, then 
    % \begin{align}\label{equ:entropy-vs-Lagrangian-b}
    %     \eta_{\alpha,Q}(\mathcal{H})
    %     = \lambda_{\alpha,Q}(\mathcal{H}).
    % \end{align}
    Let $Q$ and $\mathcal{H}$ be $r$-graphs. 
    For every $\alpha \ge 1$, 
    \begin{align*}
        \eta_{\alpha, Q}(\mathcal{H}) = \lambda_{\alpha, Q}(\mathcal{H}).
    \end{align*}
\end{proposition}

\begin{proof}[Proof of Proposition~\ref{PROP:Lagrangian-Entropy}]
    Fix $\alpha \ge 1$. 
    Let $Q$ be an $r$-graph on the vertex set $[q]$ and $\mathcal{H}$ be an $r$-graph on the vertex set $[n]$. 
    We begin by establishing the following claim.

    \begin{claim}\label{CLAIM:entropy-difference-simplify}
        Let $\mathbf{x} = (x_1, \ldots, x_n) \in \mathbb{R}^{n}$ be a nonnegative vector. 
        Let $(X_1, \ldots, X_n)$ be the random embedding of $Q$ in $\mathcal{H}$ given by the distribution
        \begin{align}\label{equ:random-Q-prob-def}
            \mathbb{P}\left((X_1, \ldots, X_n) = (i_1, \ldots, i_q)\right)
            = \frac{x_{i_1} \cdots x_{i_q}}{\beta}
            \quad\text{for every}\quad (i_1, \ldots, i_q) \in \mathrm{Inj}(Q,\mathcal{H}), 
        \end{align}
        where $\beta \coloneqq \sum_{\phi \in \mathrm{Inj}(Q,\mathcal{H})} \prod_{j \in \phi} x_j$. 
        % \begin{align*}
        %     \beta
        %     \coloneqq 
        %     \sum_{\phi \in \mathrm{Inj}(Q,\mathcal{H})} \prod_{j \in \phi} x_j.
        % \end{align*}
        %
        Also, for every $j \in [n]$, let 
        % for every $(i,j) \in [q] \times [n]$, let 
        \begin{align}
            % y_{i \to j}
            % & \coloneqq \mathbb{P}\left(X_i = j\right)
            % \coloneqq \frac{x_j}{\beta} \sum_{\phi \in \mathrm{Inj}(Q,\mathcal{H}, i \to j)} \prod_{k\in \phi - j} x_k, \quad\text{and} \label{equ:random-yij-prob-def} \\
            y_j 
            % & \coloneqq \frac{1}{q} \sum_{i \in [q]} y_{i \to j} 
            \coloneqq \frac{1}{q} \frac{x_j}{\beta} \sum_{\phi \in \mathrm{Inj}(Q,\mathcal{H}, j)} \prod_{k\in \phi - j} x_k. \label{equ:random-yj-prob-def}
        \end{align}
        %
        Then we have 
        \begin{align*}
            \mathbb{H}[(X_1, \ldots, X_{q})] - \frac{q}{\alpha} \mathbb{H}[\mu(X_1, \ldots, X_q)] 
            = \log_{2} \beta - \frac{q}{\alpha} \sum_{j \in [n]} y_j \log_{2}\left(\frac{x_j^{\alpha}}{y_j}\right). 
            % = \log_{2}\beta - \frac{1}{\alpha} \sum_{i \in [q]} \sum_{j \in [n]} y_{i \to j} \log_{2} \left(\frac{x_{j}^{\alpha}}{y_{i \to j}}\right).
            % = \log_{2}\beta - \frac{1}{\alpha} \sum_{i=1}^{q}\sum_{j=1}^{n} y_{i\to j} \log_{2} x_{j}. 
        \end{align*}
    \end{claim}
    \begin{proof}[Proof of Claim~\ref{CLAIM:entropy-difference-simplify}]
        By definition~\eqref{equ:random-Q-prob-def}, we have 
        \begin{align*}
            \mathbb{H}[(X_1, \ldots, X_{q})]
            & = \sum_{(i_1, \ldots, i_{q}) \in \mathrm{Inj}(Q,\mathcal{H})} - \frac{x_{i_1} \cdots x_{i_q}}{\beta} \log_{2} \frac{x_{i_1} \cdots x_{i_q}}{\beta} \\
            & = \sum_{(i_1, \ldots, i_{q}) \in \mathrm{Inj}(Q,\mathcal{H})} \frac{x_{i_1} \cdots x_{i_q}}{\beta} \left( \log_{2} \beta - \sum_{j\in [q]} \log_{2}x_{i_j} \right) \\
            & = \log_{2}\beta - \sum_{(i_1, \ldots, i_{q}) \in \mathrm{Inj}(Q,\mathcal{H})} \frac{x_{i_1} \cdots x_{i_q}}{\beta} \sum_{j\in [q]} \log_{2}x_{i_j} \\
            & = \log_{2}\beta - \sum_{j\in [n]} q y_{j} \log_{2} x_j. 
            % & = \log_{2}\beta - \sum_{i \in [q]}\sum_{j\in [n]} y_{i \to j} \log_{2} x_j. 
        \end{align*}
        %
        For every $(i,j) \in [q] \times [n]$, let 
        \begin{align*}
             y_{i \to j}
            \coloneqq \frac{x_j}{\beta} \sum_{\phi \in \mathrm{Inj}(Q,\mathcal{H}, i \to j)} \prod_{k\in \phi - j} x_k. 
        \end{align*}
        %
        Note that $P(X_i = j) = y_{i \to j}$ for every $(i,j) \in [q] \times [n]$ and $y_j = \sum_{i \in [q]}y_{i\to j}/q$ for every $j \in [n]$. 
        Therefore, we have 
        \begin{align*}
            \mathbb{H}[\mu(X_{1}, \ldots, X_q)]
            = - \sum_{j \in [n]} y_j \log_2 y_j, 
            % = \sum_{j \in [n]} - y_{i \to j} \log_{2} y_{i \to j}.  
        \end{align*}
        %
        and hence, 
        \begin{align*}
            & \quad \mathbb{H}\left[ (X_1, \ldots, X_{q}) \right] - \frac{q}{\alpha}  \mathbb{H}[\mu(X_{1}, \ldots, X_q)] \\
            & = \log_{2}\beta - \sum_{j\in [n]} q y_{j} \log_{2} x_j + \frac{q}{\alpha}\sum_{j \in [n]} y_j \log_2 y_j 
            % & = \log_{2}\beta - \frac{q}{\alpha} \left( \sum_{j\in [n]} y_j \left(\alpha \log_2 x_j - \log_2 y_j \right)\right) \\
            = \log_{2} \beta - \frac{q}{\alpha} \sum_{j \in [n]} y_j \log_{2}\left(\frac{x_j^{\alpha}}{y_j}\right). 
            % & = \log_{2}\beta - \sum_{i \in [q]}\sum_{j\in [n]} y_{i \to j} \log_{2} x_j + \frac{1}{\alpha}\sum_{i \in [q]} \sum_{j \in [n]} y_{i \to j} \log_{2} y_{i \to j} \\
            % & = \log_{2}\beta - \frac{1}{\alpha}\left( \alpha \sum_{i \in [q]}\sum_{j\in [n]} y_{i \to j} \log_{2} x_j + \sum_{i \in [q]} \sum_{j \in [n]} y_{i \to j} \log_{2} y_{i \to j} \right) \\
            % & = \log_{2}\beta - \frac{1}{\alpha} \sum_{i \in [q]} \sum_{j \in [n]} y_{i \to j} \log_{2} \left(\frac{x_{j}^{\alpha}}{y_{i \to j}}\right).
        \end{align*}
        This proves Claim~\ref{CLAIM:entropy-difference-simplify}. 
    \end{proof}%CLAIM

    Now we prove that $\eta_{\alpha, Q}(\mathcal{H}) \le \lambda_{\alpha, Q}(\mathcal{H})$. 
    Fix a nonnegative optimal vector $\mathbf{x} = (x_1, \ldots, x_n) \in \mathrm{OPT}_{\alpha,Q}(\mathcal{H})$, whose existence is guaranteed by Fact~\ref{FACT:nonnegative-optimal-vector}. 
    % It follows from Lemma~\ref{LEMMA:Lagrange-multiplier} that 
    % \begin{align*}
    %     \partial_{i}
    % \end{align*}
    It follows from the definition of $\mathrm{OPT}_{\alpha,Q}(\mathcal{H})$ that 
    \begin{align*}
        \beta
        \coloneqq \sum_{\phi \in \mathrm{Inj}(Q,\mathcal{H})} \prod_{j \in \phi} x_j
        = \lambda_{\alpha,Q}(\mathcal{H}). 
    \end{align*}
    %
    It follows from the convexity of the function $\log_{2} \colon \mathbb{R}_{>0} \to \mathbb{R}$ that 
    \begin{align*}
        \sum_{j \in [n]} y_j \log_{2}\left(\frac{x_j^{\alpha}}{y_j}\right) 
        \le \log_{2}\left(\sum_{j \in [n]} y_j \cdot \frac{x_j^{\alpha}}{y_j}\right)
        = \log_{2} \left(\sum_{j \in [n]} x_j^{\alpha}\right)
        = \log_{2} 1
        = 0. 
    \end{align*}
    %
    So, by Claim~\ref{CLAIM:entropy-difference-simplify}, we obtain 
    \begin{align*}
        \mathbb{H}\left[ (X_1, \ldots, X_{q}) \right] - \frac{q}{\alpha}  \mathbb{H}[\mu(X_{1}, \ldots, X_q)]
        \ge \log_{2}\beta
        = \log_{2} \lambda_{\alpha,Q}(\mathcal{H}). 
    \end{align*}
    which proves that $\eta_{\alpha, Q}(\mathcal{H}) \ge \lambda_{\alpha, Q}(\mathcal{H})$. 

    Now we consider the other direction. 
    Let $(X_1, \ldots, X_{q})$ be a random embedding of $Q$ in $\mathcal{H}$ that maximizes $\mathbb{H}[(X_1, \ldots, X_{q})] - \frac{q}{\alpha} \mathbb{H}[\mu(X_1, \ldots, X_q)]$. 
    % \begin{align*}
    %     \mathbb{H}[(X_1, \ldots, X_{q})] - \frac{q}{\alpha} \mathbb{H}[\mu(X_1, \ldots, X_q)]. 
    % \end{align*}
    %
    For every $\phi \in \mathrm{Inj}(Q, \mathcal{H})$, let 
    \begin{align*}
        p_{\phi}
        \coloneqq \mathbb{P}\left((X_1, \ldots, X_q) = \phi \right). 
    \end{align*}
    %
    We say that two embeddings $\phi, \psi \in \mathrm{Inj}(Q,\mathcal{H})$ are equivalent, denoted $\phi \sim \psi$, if there exists an automorphism $\theta \in \mathrm{Aut}(Q)$ such that $\phi \circ \theta = \psi$. 
    It follows from Property~\ref{prop:random-embedding-b} that $p_{\phi} = p_{\psi}$ whenever $\phi\sim \psi$. 
    
    Let $\mathcal{N}(Q,\mathcal{H}) \coloneqq \mathrm{Inj}(Q,\mathcal{H})/\sim$ denote the set of equivalence classes in $\mathrm{Inj}(Q,\mathcal{H})$. 
    Let $[X_1, \ldots, X_q]$ denote the probability distribution of $(X_1, \ldots, X_q)$ induced on $\mathcal{N}(Q,\mathcal{H})$, that is, each equivalence class $[\phi] \in \mathcal{N}(Q,\mathcal{H})$, 
    \begin{align}\label{equ:def-q-phi}
        q_{\phi}
        \coloneqq \mathbb{P}\left([X_1, \ldots, X_q] = [\phi] \right)
        = \sum_{\psi \in [\phi]} \mathbb{P}\left((X_1, \ldots, X_q) = \psi\right) 
        = p_{\phi} |\mathrm{Aut}(Q)|. 
    \end{align}
    %
    Then, we have  
    \begin{align}\label{equ:HX1-Xq-expression}
        \mathbb{H}\left[ (X_1, \ldots, X_q) \right]
        & = \mathbb{H}\left[ (X_1, \ldots, X_q) \mid [X_1, \ldots, X_q] \right] + \mathbb{H}\left[ [X_1, \ldots, X_q] \right] \notag \\
        & = \log_{2} |\mathrm{Aut}(Q)| - \sum_{[\phi] \in \mathcal{N}(Q,\mathcal{H})} q_{\phi} \log_{2} q_{\phi}. 
    \end{align}
    %
    For each $j \in [n]$, define  
    \begin{align*}
        \mathcal{N}(Q, \mathcal{H}, j)
        \coloneqq \big\{  [\phi] \in \mathcal{N}(Q, \mathcal{H}) \colon j \in \phi \big\}, 
    \end{align*}
    %
    and let 
    \begin{align}\label{equ:def-yj}
        y_j 
        \coloneqq \frac{1}{q} \sum_{\phi \in \mathrm{Inj}(Q,\mathcal{H}, j)}p_{\phi}
        = \frac{1}{q} \sum_{[\phi] \in \mathcal{N}(Q, \mathcal{H}, j)} q_{\phi}
        \quad\text{and}\quad 
        x_{j}
        \coloneqq  y_j^{\frac{1}{\alpha}}
        = \left(\frac{1}{q} \sum_{[\phi] \in \mathcal{N}(Q, \mathcal{H}, j)} q_{\phi}\right)^{1/\alpha}.  
    \end{align}
    %
    Note that 
    \begin{align*}
        \mathbb{P}\left( \mu(X_1, \ldots, X_q) = j \right)
        = y_j
        \quad\text{for every}\quad j \in [n],  
    \end{align*}
    %
    and thus, 
    \begin{align*}
        \mathbb{H}\left[\mu(X_1, \ldots, X_q)\right]
        = - \sum_{j\in [n]} y_j \log_{2} y_j 
        = - \sum_{j \in [n]} x_j^{\alpha} \log_{2} x_j^{\alpha}. 
    \end{align*}
    %
    Combining it with~\eqref{equ:HX1-Xq-expression}, we conclude that the distribution $\mathbf{q} \coloneqq \left\{ q_{\phi} \colon [\phi] \in \mathcal{N}(Q, \mathcal{H}) \right\}$ maximizes the objective 
    \begin{align*}
        \Phi(\mathbf{q})
        \coloneqq 
        - \sum_{[\phi] \in \mathcal{N}(Q,\mathcal{H})} q_{\phi} \log_{2} q_{\phi} + \frac{q}{\alpha} \sum_{j \in [n]} x_{j}^{\alpha} \log_{2} x_{j}^{\alpha}, 
    \end{align*}
    subject to the constraints 
    \begin{align*}
        q_{\phi} \ge 0 ~\text{for every}~ [\phi] \in \mathcal{N}(Q,\mathcal{H})
        \quad\text{and}\quad 
        \sum_{[\phi] \in \mathcal{N}(Q,\mathcal{H})} q_{\phi} = 1. 
    \end{align*}
    %
    It follows from the Lagrange multiplier method that 
    \begin{align*}
        -1 - \log_{2} q_{\phi} + \frac{1}{\alpha} \sum_{j \in \phi} \left(1+\log_{2} x_j^{\alpha}\right) 
        = -1+\frac{q}{\alpha} - \log_{2} \frac{q_{\phi}}{\prod_{j\in \phi} x_j}
    \end{align*}
    is constant for all $[\phi] \in \mathcal{N}(Q,\mathcal{H})$ with $q_{\phi} > 0$. 
    Equivalently, there exists a constant $\gamma$ such that 
    \begin{align}\label{equ:ratio-is-constant}
        \frac{q_{\phi}}{\prod_{j \in \phi} x_j} = \gamma
    \end{align}
    for all $[\phi] \in \mathcal{N}(Q,\mathcal{H})$ with $q_{\phi} > 0$.

    It follows from~\eqref{equ:def-q-phi} and~\eqref{equ:ratio-is-constant} that 
    \begin{align}\label{equ:p-phi}
        \mathbb{P}\left((X_1, \ldots, X_q) = \phi \right)
        = \frac{q_{\phi}}{|\mathrm{Aut}(Q)|}
        = \frac{\gamma}{|\mathrm{Aut}(Q)|}  \prod_{j \in \phi} x_j
        = \frac{1}{\beta} \prod_{j \in \phi} x_j,  
    \end{align}
    where $\beta \coloneqq |\mathrm{Aut}(Q)|/\gamma$. 
    
    Also, notice that 
    \begin{align*}
        \mathbb{P}\left(\mu(X_1, \ldots, X_q) = j \right) 
        = y_j 
        = \frac{1}{q}\sum_{\phi \in \mathrm{Inj}(Q,\mathcal{H}, j)}p_{\phi}
        = \frac{x_j}{q \beta}  \sum_{\phi \in \mathrm{Inj}(Q,\mathcal{H}, j)} \prod_{k \in \phi - j} x_k.
        % = \frac{\gamma x_j}{q |\mathrm{Aut}(Q)|}  \sum_{\phi \in \mathrm{Inj}(Q,\mathcal{H}, j)} \prod_{k \in \phi - j} y_k. 
    \end{align*}
    %
    So it follows from Claim~\ref{CLAIM:entropy-difference-simplify} and definition~\eqref{equ:def-yj} that 
    \begin{align}\label{equ:entropy-Q-lower-bound}
        \mathbb{H}[(X_1, \ldots, X_{q})] - \frac{q}{\alpha} \mathbb{H}[\mu(X_1, \ldots, X_q)] 
        & = \log_{2} \beta - \frac{q}{\alpha} \sum_{j \in [n]} y_j \log_{2}\left(\frac{x_j^{\alpha}}{y_j}\right) \notag \\
        & = \log_{2} \beta - \frac{q}{\alpha} \sum_{j \in [n]} y_j \log_{2} 1
        = \log_{2} \beta.  
    \end{align}
    %
    Note from~\eqref{equ:p-phi} that the partition function $\beta$ satisfies 
    \begin{align*}
        \beta 
        = \sum_{\phi \in \mathrm{Inj}(Q, \mathcal{H})} \prod_{j \in \phi} x_j. 
    \end{align*}
    %
    Also observe that  
    \begin{align*}
        \sum_{i\in [n]} x_i^{\alpha} = \sum_{i \in [n]} y_i = 1
        \quad\text{and}\quad 
        x_i \ge 0 ~\text{for every}~ i \in [n], 
    \end{align*}
    which imply that $(x_1, \ldots, x_n) \in \Delta_{\alpha}^{n-1}$. 
    So, it follows from the definition of $\lambda_{\alpha, Q}(\mathcal{H})$ that $\beta \le \lambda_{\alpha, Q}(\mathcal{H})$. 
    Combining this with~\eqref{equ:entropy-Q-lower-bound}, we conclude that $\eta_{\alpha, Q}(\mathcal{H}) \le \lambda_{\alpha, Q}(\mathcal{H})$. 
    This completes the proof of Proposition~\ref{PROP:Lagrangian-Entropy}. 
\end{proof}%PROP

\end{document}
